% ============================================================
% Chapter 1: Introduction
%   Section: Objectives
%     Subsection: Specific Objectives
% Included from chapters/01-introduction/04-objectives/objectives.tex
% One per deliverable of docs/roadmap.md. Whether each was achieved is
% Section~\ref{sec:objectives-achieved}.
% ============================================================

\subsection{Specific Objectives}

\begin{description}[leftmargin=1.2cm,labelwidth=1cm]
    \item[O1] \textbf{Discovery and selection.} A flow that reads the
    configured English and Spanish sources for up to three topic
    groups, keeps what fits the publication's mission, ranks the
    candidates, and lets a person or a language model choose at most
    twenty to analyse (RQ1, RQ3).
    \item[O2] \textbf{Extraction and enrichment.} A pipeline that turns
    any article URL into clean text with its metadata, and annotates it
    with topics, entities, keywords, sentiment, quality scores,
    candidate claims and an embedding, stored in a data lake that keeps
    the lineage of every record.
    \item[O3] \textbf{Open-web fact-checking with auditable verdicts.}
    A pipeline that admits an article, selects its checkable claims,
    retrieves evidence from the open web, and returns for each claim a
    verdict that cites its sources and can be traced stage by stage
    (RQ2).
    \item[O4] \textbf{A validation set and an evaluation harness.} A
    hand-labelled set of claims from the system's own domain, and a
    harness that runs the pipeline over it and over X-Fact and reports
    every metric with its uncertainty (RQ2, RQ3, RQ4).
    \item[O5] \textbf{Model benchmarks.} A comparison of language
    models as verifiers on the labelled sets, and of a local and a
    hosted model as the editor behind the text corrector, with their
    cost and latency (RQ2).
    \item[O6] \textbf{A reader view.} A public view of the published
    articles that shows each claim's verdict, its evidence and why each
    source was trusted.
\end{description}
